\section{Relative abundance for projective analytic morphisms}
\label{sec:analytic-abundance}

The complex rational case of Theorem~\ref{thm:main} supplies the algebraic
fiber input in Fujino's relative analytic reduction
\cite[Theorem~1.10 and its proof]{FujinoRelativeAnalytic2025}.
We record only its normal rational case. Throughout this section, complex
analytic spaces are Hausdorff and second-countable, and a complex variety
means a reduced and irreducible complex analytic space, as in
\cite[p.~5]{FujinoRelativeAnalytic2025}.

Here a projective analytic morphism is proper and possesses a relatively
ample line bundle. Rational Cartierness is local on the source; a single
global Cartier index is not part of that definition. A divisor is
\(\pi\)-nef over \(Y\) if it has nonnegative intersection with every
projective integral curve contracted by \(\pi\), over every point of \(Y\).
These are the conventions of
\cite[Definitions~2.28, 2.32 and~2.48]{FujinoAnalyticMMP2022}.

\begin{corollary}[Relative analytic abundance]
\label{cor:analytic-abundance}
Let \(\pi:X\to Y\) be a projective surjective morphism of normal complex
analytic varieties. Let \((X,\Delta)\) be a log canonical pair with
\(\Delta\) an effective rational divisor and \(D=K_X+\Delta\)
\(\Q\)-Cartier. Assume that \(D\) is \(\pi\)-nef over all of \(Y\).
For every compact subset \(W\subset Y\), there exist an analytic open
neighborhood \(U\supset W\) and a positive integer \(m\) such that
\(mD|_{X_U}\) is Cartier, where \(X_U=\pi^{-1}(U)\), and the line bundle
\[
 L=\OO_{X_U}(mD|_{X_U})
\]
is generated relative to \(\pi_U:X_U\to U\). Equivalently, the evaluation
map
\[
 \pi_U^*(\pi_U)_*L\longrightarrow L
\]
is surjective.
\end{corollary}
\begin{proof}
Since \(\pi\) is proper, \(\pi^{-1}(W)\) is compact. Local rational
Cartier indices therefore have a common multiple \(m_0\) on an open
neighborhood \(O\) of \(\pi^{-1}(W)\). The image of the closed set
\(X\setminus O\) is closed and disjoint from \(W\), so there is an
open neighborhood \(V\supset W\) with \(\pi^{-1}(V)\subset O\).
Thus \(m_0D\) is Cartier over \(V\).

Fix \(P\in W\). Following
\cite[the proof of Theorem~1.10, p.~30]{FujinoRelativeAnalytic2025},
work over a sufficiently small neighborhood of \(P\) contained in \(V\)
and take a crepant dlt blow-up
\(p:(X',\Delta')\to(X,\Delta)\). Put \(\pi'=\pi\circ p\) and
\(D'=K_{X'}+\Delta'=p^*D\). For every log canonical stratum \(S\)
of \((X',\Delta')\), including \(X'\) itself, adjunction gives an
effective rational boundary \(\Delta_S\) with
\[
 K_S+\Delta_S=D'|_S.
\]
The strata are normal, and this adjoint is nef over \(\pi'(S)\).
On each normal projective component \(F\) of an analytically sufficiently
general fiber of \(S\to\pi'(S)\), restriction gives a log canonical
pair \((F,\Delta_F)\) with effective rational boundary and nef
adjoint \(K_F+\Delta_F=(K_S+\Delta_S)|_F\).
Theorem~\ref{thm:main} makes this adjoint semiample, so its Iitaka
and numerical dimensions agree. Thus \(D'|_S\) is abundant over
\(\pi'(S)\) for every stratum: this is precisely
\(\pi'\)-log abundance in
\cite[Definitions~2.10 and~2.13]{FujinoRelativeAnalytic2025}.

The nef-and-log-abundant theorem
\cite[Theorem~1.4]{FujinoRelativeAnalytic2025} now gives a relatively
generated multiple of \(D'\) over a neighborhood \(U_P\subset V\)
of \(P\). Enlarge its integer to a multiple \(m_P\) of \(m_0\);
tensor powers preserve relative generation. Normality gives
\(p_*\OO_{X'}=\OO_X\), and projection formula identifies the direct
images over \(U_P\) of \(\OO_X(m_PD)\) and
\(\OO_{X'}(m_PD')\). Their evaluation maps are related by pullback.
Generation upstairs therefore implies generation downstairs: a base point
downstairs would make every pulled-back section vanish on its nonempty
fiber.

Choose finitely many of these neighborhoods \(U_{P_1},\ldots,U_{P_s}\)
covering \(W\), and let \(m\) be a common multiple of
\(m_0,m_{P_1},\ldots,m_{P_s}\). Put \(U=\bigcup_i U_{P_i}\).
The divisor \(mD\) is Cartier over \(U\), and its line bundle is
relatively generated on each \(U_{P_i}\). Surjectivity of evaluation
is local on the base, so it holds over \(U\), as asserted.
\end{proof}

The base \(Y\) need not be algebraic, compact, or Stein, and \(W\) need
not be a Stein compact subset. The integer \(m\) may depend on the pair,
the morphism, and \(W\); no single Cartier index or relatively generated
multiple over all of a noncompact base is asserted. The nefness premise is
over all of \(Y\), not only over \(W\). Projectivity cannot be replaced here
by mere properness or by a K\"ahler hypothesis: when \(Y\) is a point it
requires \(X\) to be projective. The statement does not include real
boundaries or semi-log-canonical analytic pairs.
